\phantomsection
\section*{Введение}
\addcontentsline{toc}{section}{Введение}

\textbf{Актуальность темы.} Рынок образовательных услуг в области иностранных
языков демонстрирует устойчивый рост: всё большее число частных языковых центров
сталкивается с необходимостью автоматизации административных процессов - записи
студентов на занятия, формирования расписания и контроля успеваемости. Ручное
ведение журналов и электронных таблиц приводит к ошибкам при записи, трудностям
в отслеживании прогресса студентов и неэффективному использованию рабочего времени
преподавателей. Разработка специализированного веб-приложения, объединяющего
систему записи, расписание и модуль успеваемости в едином интерфейсе с разграничением
доступа по ролям, является актуальной практической задачей.

\textbf{Степень разработанности проблемы.} Существующие решения - Moodle,
Google Classroom, iSpring Learn - решают задачу частично: универсальные LMS
не обеспечивают запись на конкретные занятия с контролем вместимости и встроенного
управления расписанием, а CRM-системы (Bitrix24, AmoCRM) не предоставляют
функциональности для работы с учебными материалами и отслеживания успеваемости.
Специализированного открытого решения на стеке Flask + PostgreSQL, адаптированного
для небольших языковых центров, в открытом доступе не представлено.

\textbf{Объект исследования} - процессы организации учебного процесса в центре
изучения иностранных языков: запись студентов на занятия, формирование расписания
и контроль успеваемости.

\textbf{Предмет исследования} - программные средства и архитектурные решения
для реализации веб-приложения с ролевой моделью, системой записи на занятия
и агрегацией данных об успеваемости студентов.

\textbf{Целью данной работы} является проектирование и реализация веб-приложения
для центра изучения иностранных языков, обеспечивающего запись студентов на занятия,
управление расписанием, загрузку учебных материалов и отслеживание успеваемости
в рамках единой системы с разграничением доступа по ролям администратора,
преподавателя и студента.
